\documentclass[12pt,a4paper,oneside]{report}

\usepackage[utf8]{inputenc}
\usepackage[T1]{fontenc}
\usepackage[french]{babel}
\usepackage{lmodern}
\usepackage{microtype}
\usepackage{geometry}
\usepackage{xcolor}
\usepackage{graphicx}
\usepackage{booktabs}
\usepackage{tabularx}
\usepackage{longtable}
\usepackage{array}
\usepackage{enumitem}
\usepackage{titlesec}
\usepackage{fancyhdr}
\usepackage{hyperref}
\usepackage{tikz}
\usepackage{listings}
\usepackage{caption}
\usepackage{subcaption}
\usepackage{float}
\usepackage{pdflscape}
\usepackage{pifont}

\usetikzlibrary{
  arrows.meta,
  backgrounds,
  calc,
  fit,
  positioning,
  shapes.geometric,
  shapes.misc
}

\geometry{margin=2.4cm}
\setlength{\parskip}{0.55em}
\setlength{\parindent}{0pt}
\setlist[itemize]{topsep=0.35em,itemsep=0.25em}
\setlist[enumerate]{topsep=0.35em,itemsep=0.25em}

\definecolor{DevPilotBlue}{HTML}{1E3A8A}
\definecolor{DevPilotTeal}{HTML}{0F766E}
\definecolor{DevPilotGreen}{HTML}{15803D}
\definecolor{DevPilotOrange}{HTML}{C2410C}
\definecolor{DevPilotRed}{HTML}{B91C1C}
\definecolor{DevPilotGray}{HTML}{334155}
\definecolor{LightBlue}{HTML}{EFF6FF}
\definecolor{LightTeal}{HTML}{ECFDF5}
\definecolor{LightOrange}{HTML}{FFF7ED}
\definecolor{LightGray}{HTML}{F8FAFC}
\definecolor{CodeBg}{HTML}{F1F5F9}

\hypersetup{
  colorlinks=true,
  linkcolor=DevPilotBlue,
  citecolor=DevPilotTeal,
  urlcolor=DevPilotBlue,
  pdftitle={Rapport de projet DevPilot AI},
  pdfauthor={DevPilot AI}
}

\pagestyle{fancy}
\fancyhf{}
\lhead{\nouppercase{\leftmark}}
\rhead{DevPilot AI}
\cfoot{\thepage}
\renewcommand{\headrulewidth}{0.4pt}

\titleformat{\chapter}
  {\Huge\bfseries\color{DevPilotBlue}}
  {\thechapter}
  {1em}
  {}
\titleformat{\section}
  {\Large\bfseries\color{DevPilotGray}}
  {\thesection}
  {0.75em}
  {}
\titleformat{\subsection}
  {\large\bfseries\color{DevPilotTeal}}
  {\thesubsection}
  {0.75em}
  {}

\lstdefinestyle{devpilot}{
  backgroundcolor=\color{CodeBg},
  basicstyle=\ttfamily\small,
  breaklines=true,
  columns=fullflexible,
  frame=single,
  rulecolor=\color{DevPilotGray},
  keywordstyle=\color{DevPilotBlue}\bfseries,
  commentstyle=\color{DevPilotGreen},
  stringstyle=\color{DevPilotOrange},
  showstringspaces=false
}
\lstset{style=devpilot}

\newcommand{\cmark}{\ding{51}}
\newcommand{\xmark}{\ding{55}}
\newcommand{\placeholder}[2]{
  \begin{figure}[H]
    \centering
    \fcolorbox{DevPilotGray}{LightGray}{
      \begin{minipage}[c][#2][c]{0.88\textwidth}
        \centering
        \textbf{Placeholder capture d'écran}\\[0.4em]
        #1
      \end{minipage}
    }
    \caption{#1}
  \end{figure}
}

\newcolumntype{Y}{>{\raggedright\arraybackslash}X}
\newcolumntype{L}[1]{>{\raggedright\arraybackslash}p{#1}}

\begin{document}

\begin{titlepage}
  \centering
  \vspace*{1cm}
  {\Large Rapport de projet d'ingénierie\par}
  \vspace{1.5cm}
  {\Huge\bfseries\color{DevPilotBlue} DevPilot AI\par}
  \vspace{0.5cm}
  {\Large Plateforme agentique de RAG privé pour l'intelligence documentaire logicielle\par}
  \vspace{1.2cm}
  \begin{tikzpicture}
    \node[
      rounded corners=8pt,
      draw=DevPilotBlue,
      very thick,
      fill=LightBlue,
      inner sep=18pt,
      text width=0.82\textwidth,
      align=center
    ] {
      \Large\bfseries Architecture backend moderne, pipelines asynchrones,\\
      retrieval hybride, évaluation de réponses, observabilité et préparation production
    };
  \end{tikzpicture}
  \vfill
  \begin{tabular}{rl}
    \textbf{Projet :} & DevPilot AI \\
    \textbf{Périmètre :} & Phases réalisées jusqu'à la consolidation avancée \\
    \textbf{Encadrement :} & Rapport académique et technique \\
    \textbf{Date :} & \today
  \end{tabular}
  \vfill
  {\large Année universitaire 2025--2026\par}
\end{titlepage}

\pagenumbering{roman}

\chapter*{Résumé}
\addcontentsline{toc}{chapter}{Résumé}
DevPilot AI est une plateforme d'intelligence documentaire privée destinée à assister l'analyse de projets logiciels. Le système permet à un utilisateur authentifié de travailler dans des espaces isolés, d'importer des documents, de déclencher leur traitement asynchrone, d'indexer leurs fragments dans une base vectorielle, puis d'interroger ces connaissances à travers un workflow de Retrieval-Augmented Generation agentique.

Le projet dépasse le cadre d'un simple prototype conversationnel. Il combine une architecture backend FastAPI modulaire, une interface Next.js, une persistance PostgreSQL, une file Redis/Celery, une base vectorielle Qdrant, des fournisseurs IA extensibles, une stratégie de retrieval hybride, du reranking, une évaluation automatique de la qualité des réponses et une couche d'observabilité fondée sur des logs structurés, Prometheus et Grafana.

Les phases de consolidation avancée ont renforcé la qualité logicielle, la stabilité des services, la sécurité des échanges, la validation des API, la robustesse du pipeline d'ingestion et la préparation production-ready. Le rapport présente la problématique, l'étude de l'existant, l'analyse des besoins, la conception, le system design, les choix techniques, la sécurité, les tests, l'observabilité, le DevOps, les résultats, les limites et les perspectives.

\textbf{Mots-clés :} RAG, FastAPI, Next.js, PostgreSQL, Qdrant, Celery, Redis, Ollama, observabilité, architecture logicielle, sécurité, DevOps.

\chapter*{Abstract}
\addcontentsline{toc}{chapter}{Abstract}
DevPilot AI is a private document intelligence platform designed to support software project analysis. The system allows authenticated users to work in isolated workspaces, upload documents, process them asynchronously, index chunks in a vector database, and query this knowledge through an agentic Retrieval-Augmented Generation workflow.

The project is not limited to a basic chatbot prototype. It combines a modular FastAPI backend, a Next.js frontend, PostgreSQL persistence, Redis/Celery asynchronous processing, Qdrant vector search, extensible AI providers, hybrid retrieval, reranking, automatic answer evaluation, and an observability layer based on structured logs, Prometheus, and Grafana.

The advanced consolidation phases improved software quality, service stability, API validation, secure exchanges, ingestion robustness, monitoring, and production readiness. This report presents the problem statement, state of the art, requirements analysis, design, system design, technical decisions, security, tests, observability, DevOps, results, limitations, and perspectives.

\textbf{Keywords:} RAG, FastAPI, Next.js, PostgreSQL, Qdrant, Celery, Redis, Ollama, observability, software architecture, security, DevOps.

\tableofcontents
\listoffigures
\listoftables
\clearpage
\pagenumbering{arabic}

\chapter*{Plan détaillé du rapport}
\addcontentsline{toc}{chapter}{Plan détaillé du rapport}
Le rapport est structuré selon une logique d'ingénierie complète : partir du besoin, formaliser la conception, justifier l'architecture, prouver la qualité par les tests et analyser la préparation à l'exploitation.

\begin{enumerate}
  \item \textbf{Introduction générale} : contexte, problématique, objectifs, périmètre et démarche.
  \item \textbf{Étude de l'existant} : limites des assistants IA génériques, plateformes RAG, outils DevOps et observabilité.
  \item \textbf{Analyse des besoins} : acteurs, cas d'utilisation, exigences fonctionnelles et non fonctionnelles.
  \item \textbf{Analyse et conception} : modèle de domaine, règles métier, séparation des responsabilités.
  \item \textbf{System Design} : architecture logique, déploiement, flux d'ingestion, flux RAG, scalabilité et résilience.
  \item \textbf{Réalisation technique} : backend, frontend, services IA, workers, persistance et configuration.
  \item \textbf{Sécurité} : authentification, isolation multi-workspace, validation, secrets et risques.
  \item \textbf{Tests et qualité logicielle} : stratégie de validation, résultats, critères d'acceptation et dette connue.
  \item \textbf{Monitoring et observabilité} : logs structurés, métriques, Prometheus, Grafana et corrélation.
  \item \textbf{DevOps et exploitation} : Docker Compose, Makefile, migrations, environnements et préparation production.
  \item \textbf{Résultats, limites et perspectives} : bilan technique, limites réelles et axes d'amélioration.
  \item \textbf{Conclusion, bibliographie et annexes} : synthèse, références, API, modèle de données et commandes.
\end{enumerate}

\chapter{Introduction générale}

\section{Contexte}
Les équipes de développement produisent une quantité croissante de documents techniques : spécifications, journaux de décision, README, plans de test, tickets, extraits de code, notes d'architecture et rapports de validation. Ces documents contiennent une connaissance importante, mais celle-ci reste souvent dispersée, difficile à interroger et insuffisamment reliée aux workflows de développement.

L'émergence des grands modèles de langage a rendu possible l'assistance intelligente à la lecture et à la synthèse. Toutefois, un modèle généraliste ne connaît pas les documents privés d'un projet et peut produire des réponses plausibles mais non vérifiées. Le besoin réel n'est donc pas seulement de générer du texte, mais de construire une plateforme capable de retrouver les sources pertinentes, de citer les éléments utilisés, de contrôler la qualité des réponses et de s'intégrer dans une architecture exploitable.

DevPilot AI répond à ce besoin par une plateforme privée de RAG agentique. Le système est conçu comme une application complète : backend, frontend, base relationnelle, base vectorielle, traitement asynchrone, agents spécialisés, tests automatisés, observabilité et environnement conteneurisé.

\section{Problématique}
La problématique centrale du projet peut être formulée ainsi :

\begin{quote}
Comment concevoir une plateforme d'intelligence documentaire pour projets logiciels qui soit à la fois fiable, sécurisée, traçable, extensible et suffisamment proche d'un environnement de production pour démontrer une réelle maturité d'ingénierie ?
\end{quote}

Cette problématique implique plusieurs sous-questions :
\begin{itemize}
  \item comment isoler les données privées de plusieurs utilisateurs ou espaces de travail ;
  \item comment transformer des documents hétérogènes en fragments exploitables par un moteur de recherche vectoriel ;
  \item comment combiner recherche sémantique, recherche lexicale, reranking et génération contrôlée ;
  \item comment mesurer la qualité des réponses et réduire le risque d'hallucination ;
  \item comment rendre le système observable, testable et maintenable.
\end{itemize}

\section{Objectifs du projet}
Les objectifs principaux sont les suivants :
\begin{itemize}
  \item fournir une expérience web permettant l'authentification, la gestion de workspaces, l'import de documents et le chat documentaire ;
  \item mettre en place un backend modulaire exposant des API REST sécurisées ;
  \item traiter les documents de manière asynchrone afin de ne pas bloquer l'expérience utilisateur ;
  \item indexer les fragments dans Qdrant avec des embeddings générés localement via Ollama ;
  \item interroger les documents via un workflow RAG agentique incluant routing, réécriture, retrieval, reranking, génération, évaluation et correction ;
  \item enregistrer les conversations, citations, usages LLM, évaluations et traces d'agents ;
  \item instrumenter le système avec des métriques et des logs structurés ;
  \item fournir une base de tests et de validation compatible avec une démarche production-ready.
\end{itemize}

\section{Périmètre}
Le rapport couvre le projet jusqu'aux phases de consolidation avancée qui renforcent la robustesse, la qualité logicielle, l'observabilité et la préparation à l'exploitation. Il ne présente pas de fonctionnalités situées en dehors de ce périmètre.

\section{Démarche adoptée}
La démarche suivie est incrémentale. Les fondations techniques sont d'abord posées : monorepo, Docker Compose, backend, frontend, base de données et configuration. Les fonctionnalités métier sont ensuite ajoutées progressivement : authentification, workspaces, documents, ingestion, embeddings, retrieval, chat, évaluation et agents. Enfin, les dernières étapes du périmètre consolident le système par des tests plus complets, des métriques, des logs structurés, des tableaux de bord et des commandes d'exploitation.

\chapter{Étude de l'existant}

\section{Assistants IA généralistes}
Les assistants IA généralistes offrent une interaction naturelle avec l'utilisateur, mais ils sont limités lorsqu'il s'agit de documents privés ou de connaissances projet. Leur contexte est borné par la fenêtre de prompt, leurs réponses ne sont pas toujours sourcées et ils ne possèdent pas nativement de mécanisme d'isolation par workspace.

\section{Recherche documentaire classique}
Les moteurs de recherche classiques reposent surtout sur la correspondance lexicale. Ils sont efficaces pour retrouver des mots exacts, mais moins adaptés aux questions exprimées avec des formulations différentes. Dans un contexte logiciel, une requête comme \og traitement asynchrone des fichiers \fg{} doit pouvoir retrouver des passages mentionnant Celery, Redis, queue, worker ou indexing, même si les mots de la question ne sont pas tous présents.

\section{Retrieval-Augmented Generation}
Le RAG combine un système de recherche et un modèle génératif. Le modèle ne répond pas uniquement à partir de ses connaissances internes : il reçoit un contexte extrait de documents indexés. Cette approche réduit les hallucinations et améliore la traçabilité, à condition de maîtriser la qualité du retrieval, la taille du contexte, le prompt, les citations et l'évaluation.

\section{Plateformes et briques comparables}
\begin{table}[H]
  \centering
  \caption{Comparaison synthétique avec les familles d'outils existantes}
  \begin{tabularx}{\textwidth}{L{3.2cm}Y Y}
    \toprule
    \textbf{Famille} & \textbf{Apports} & \textbf{Limites pour DevPilot AI} \\
    \midrule
    Chatbots généralistes & Interaction naturelle, génération rapide, faible barrière d'entrée & Faible contrôle sur les données privées, réponses rarement reliées à des citations internes, observabilité applicative limitée \\
    Moteurs documentaires & Indexation robuste, recherche textuelle, filtrage & Compréhension sémantique limitée, pas de génération contextualisée \\
    Frameworks RAG & Accélèrent la construction de pipelines IA & Peuvent masquer l'architecture, les responsabilités et les mécanismes de persistance \\
    Outils DevOps & Automatisation, monitoring, déploiement & Ne répondent pas seuls au besoin de compréhension documentaire intelligente \\
    \bottomrule
  \end{tabularx}
\end{table}

\section{Positionnement de DevPilot AI}
DevPilot AI se positionne comme une plateforme intégrée orientée ingénierie. Le projet ne cherche pas seulement à appeler un modèle de langage, mais à concevoir un système complet où chaque réponse peut être reliée à des documents, des chunks, des stratégies de recherche, des agents, des métriques et des traces persistées.

\chapter{Analyse des besoins}

\section{Acteurs}
\begin{table}[H]
  \centering
  \caption{Acteurs du système}
  \begin{tabularx}{\textwidth}{L{3cm}Y}
    \toprule
    \textbf{Acteur} & \textbf{Rôle} \\
    \midrule
    Utilisateur authentifié & Crée ou consulte des workspaces, importe des documents, pose des questions, consulte les réponses et citations. \\
    Administrateur & Dispose d'un niveau d'accès supérieur pour la supervision et les opérations. \\
    Worker Celery & Exécute les traitements asynchrones d'extraction, nettoyage, chunking, embedding et indexing. \\
    Fournisseur IA & Génère des embeddings, produit les réponses et peut évaluer ou reranker selon la configuration. \\
    Système d'observabilité & Collecte métriques, logs et signaux de santé pour diagnostiquer le système. \\
    \bottomrule
  \end{tabularx}
\end{table}

\section{Besoins fonctionnels}
\begin{longtable}{L{2.2cm}L{4.3cm}Y}
  \caption{Besoins fonctionnels principaux}\\
  \toprule
  \textbf{Code} & \textbf{Besoin} & \textbf{Description} \\
  \midrule
  \endfirsthead
  \toprule
  \textbf{Code} & \textbf{Besoin} & \textbf{Description} \\
  \midrule
  \endhead
  BF-01 & Authentification & Inscription, connexion, émission de JWT, récupération du profil courant. \\
  BF-02 & Workspaces & Création, consultation, liste et isolation des espaces de travail. \\
  BF-03 & Import de documents & Upload de fichiers PDF, TXT et Markdown avec validation de type et taille. \\
  BF-04 & Traitement asynchrone & Passage des documents par les états queued, processing, indexed ou failed. \\
  BF-05 & Chunking & Extraction, nettoyage, découpage et enrichissement des fragments. \\
  BF-06 & Indexation vectorielle & Génération d'embeddings et stockage dans Qdrant. \\
  BF-07 & Retrieval & Recherche sémantique, lexicale et hybride dans le périmètre du workspace. \\
  BF-08 & Chat RAG & Génération de réponses contextualisées avec citations. \\
  BF-09 & Workflow agentique & Routing, réécriture, retrieval, reranking, génération, évaluation et correction. \\
  BF-10 & Évaluation & Calcul de scores de faithfulness, relevance, context precision et hallucination. \\
  BF-11 & Historique & Persistance des conversations, messages, citations, évaluations et usages LLM. \\
  BF-12 & Monitoring & Exposition de métriques HTTP, ingestion, retrieval, RAG, évaluation et tokens. \\
  \bottomrule
\end{longtable}

\section{Besoins non fonctionnels}
\begin{table}[H]
  \centering
  \caption{Exigences non fonctionnelles}
  \begin{tabularx}{\textwidth}{L{3.2cm}Y}
    \toprule
    \textbf{Exigence} & \textbf{Application dans DevPilot AI} \\
    \midrule
    Sécurité & JWT, hash bcrypt, CORS contrôlé, isolation par workspace, non-exposition des secrets. \\
    Maintenabilité & Découpage en routes, schémas, services, modèles, workers, providers et agents. \\
    Observabilité & Logs JSON corrélés par request id, métriques Prometheus, dashboard Grafana. \\
    Robustesse & Retries Celery, gestion d'erreurs applicatives, fallback heuristique pour l'évaluation et le reranking. \\
    Extensibilité & Abstraction des providers LLM et embedding, support Ollama et points d'extension Gemini/OpenAI. \\
    Performance & Traitement asynchrone, recherche vectorielle, batch embeddings, limitation du contexte RAG. \\
    Portabilité & Docker Compose, variables d'environnement, Makefile, migrations Alembic. \\
    Testabilité & Tests unitaires backend, tests d'intégration live, scénarios ingestion/RAG et smoke tests frontend. \\
    \bottomrule
  \end{tabularx}
\end{table}

\section{Cas d'utilisation}
\begin{figure}[H]
  \centering
  \begin{tikzpicture}[
    actor/.style={draw, rounded corners=2pt, fill=LightBlue, minimum width=2.9cm, minimum height=0.9cm},
    usecase/.style={draw, ellipse, fill=LightTeal, minimum width=3.8cm, minimum height=1cm, align=center},
    line/.style={-Latex, thick, color=DevPilotGray}
  ]
    \node[actor] (user) {Utilisateur};
    \node[actor, below=1.2cm of user] (admin) {Administrateur};

    \node[usecase, right=3.2cm of user] (auth) {S'authentifier};
    \node[usecase, below=0.85cm of auth] (ws) {Gérer les workspaces};
    \node[usecase, below=0.85cm of ws] (doc) {Importer des documents};
    \node[usecase, below=0.85cm of doc] (chat) {Interroger le RAG};
    \node[usecase, below=0.85cm of chat] (obs) {Consulter l'état système};

    \draw[line] (user) -- (auth);
    \draw[line] (user) -- (ws);
    \draw[line] (user) -- (doc);
    \draw[line] (user) -- (chat);
    \draw[line] (admin) -- (obs);
    \draw[line] (admin) -- (ws);
  \end{tikzpicture}
  \caption{Cas d'utilisation principaux de DevPilot AI}
\end{figure}

\chapter{Analyse et conception}

\section{Principes de conception}
La conception de DevPilot AI repose sur cinq principes :
\begin{itemize}
  \item \textbf{séparation des responsabilités} : les routes HTTP ne contiennent pas la logique métier, qui est portée par les services ;
  \item \textbf{isolation métier} : toutes les opérations documentaires et conversationnelles sont rattachées à un workspace ;
  \item \textbf{traitement asynchrone} : les opérations lourdes sont exécutées par Celery ;
  \item \textbf{traçabilité} : les réponses sont reliées aux chunks, aux messages, aux agents et aux évaluations ;
  \item \textbf{extensibilité} : les providers IA et les stratégies de retrieval peuvent évoluer sans réécrire l'application.
\end{itemize}

\section{Découpage applicatif}
Le backend est organisé en modules cohérents :
\begin{itemize}
  \item \texttt{api} : routes FastAPI et dépendances ;
  \item \texttt{core} : configuration, sécurité, exceptions, logs, métriques et observabilité ;
  \item \texttt{db} : sessions SQLAlchemy et base déclarative ;
  \item \texttt{models} : entités relationnelles ;
  \item \texttt{schemas} : contrats Pydantic ;
  \item \texttt{services} : logique métier et cas d'utilisation ;
  \item \texttt{workers} : configuration Celery et tâches asynchrones ;
  \item \texttt{agents} : workflow RAG agentique ;
  \item \texttt{providers} : abstraction des fournisseurs IA ;
  \item \texttt{utils} : nettoyage et découpage de texte.
\end{itemize}

\section{Modèle de domaine}
\begin{figure}[H]
  \centering
  \resizebox{\textwidth}{!}{
  \begin{tikzpicture}[
    entity/.style={draw, rounded corners=4pt, fill=LightGray, minimum width=3cm, align=left, font=\small},
    rel/.style={-Latex, thick, color=DevPilotGray}
  ]
    \node[entity] (user) {\textbf{User}\\id\\email\\role\\is\_active};
    \node[entity, right=2.1cm of user] (workspace) {\textbf{Workspace}\\id\\name\\owner\_id};
    \node[entity, below=1.2cm of workspace] (member) {\textbf{WorkspaceMember}\\workspace\_id\\user\_id\\role};
    \node[entity, right=2.1cm of workspace] (document) {\textbf{Document}\\filename\\status\\storage\_path};
    \node[entity, right=2.1cm of document] (chunk) {\textbf{Chunk}\\content\\chunk\_index\\vector\_id};
    \node[entity, below=1.4cm of document] (conversation) {\textbf{Conversation}\\title\\workspace\_id};
    \node[entity, right=2.1cm of conversation] (message) {\textbf{Message}\\role\\content};
    \node[entity, below=1.3cm of message] (retrieved) {\textbf{RetrievedChunk}\\score\\rank\\strategy};
    \node[entity, right=2.1cm of message] (eval) {\textbf{Evaluation}\\faithfulness\\relevance\\hallucination};
    \node[entity, above=1.2cm of eval] (agent) {\textbf{AgentRun}\\agent\_type\\status\\latency};
    \node[entity, below=1.2cm of eval] (usage) {\textbf{LLMUsage}\\provider\\model\\tokens};

    \draw[rel] (user) -- node[above]{owns} (workspace);
    \draw[rel] (user) -- (member);
    \draw[rel] (workspace) -- (member);
    \draw[rel] (workspace) -- (document);
    \draw[rel] (document) -- (chunk);
    \draw[rel] (workspace) -- (conversation);
    \draw[rel] (conversation) -- (message);
    \draw[rel] (message) -- (retrieved);
    \draw[rel] (retrieved) -- (chunk);
    \draw[rel] (message) -- (eval);
    \draw[rel] (message) -- (agent);
    \draw[rel] (message) -- (usage);
  \end{tikzpicture}}
  \caption{Modèle conceptuel simplifié du domaine}
\end{figure}

\section{Règles métier structurantes}
\begin{itemize}
  \item Un utilisateur doit être authentifié pour accéder aux fonctionnalités métier.
  \item Un document appartient toujours à un workspace.
  \item Un utilisateur non membre ne peut pas lire, importer, rechercher ou interroger les données d'un autre workspace.
  \item Un document supprimé logiquement ne doit plus apparaître dans les listes actives.
  \item Les chunks sont remplacés lors d'un nouveau traitement du même document afin d'éviter les index obsolètes.
  \item Une réponse RAG doit conserver ses citations, son évaluation et ses traces d'agents.
  \item Les secrets ne doivent jamais être renvoyés par les endpoints de configuration.
\end{itemize}

\chapter{System Design}

\section{Vue d'ensemble}
Le system design de DevPilot AI combine une architecture web classique et une chaîne IA spécialisée. L'application est composée d'un frontend Next.js, d'un backend FastAPI, d'une base relationnelle PostgreSQL, d'un broker Redis, d'un worker Celery, d'une base vectorielle Qdrant, d'un runtime IA local Ollama et d'une couche Prometheus/Grafana.

\begin{figure}[H]
  \centering
  \resizebox{\textwidth}{!}{
  \begin{tikzpicture}[
    comp/.style={draw, rounded corners=6pt, fill=LightBlue, minimum width=3.2cm, minimum height=1cm, align=center},
    data/.style={draw, cylinder, shape border rotate=90, aspect=0.25, fill=LightTeal, minimum width=2.6cm, minimum height=1.1cm, align=center},
    infra/.style={draw, rounded corners=6pt, fill=LightOrange, minimum width=3.2cm, minimum height=1cm, align=center},
    obs/.style={draw, rounded corners=6pt, fill=LightGray, minimum width=3.2cm, minimum height=1cm, align=center},
    line/.style={-Latex, thick, color=DevPilotGray}
  ]
    \node[comp] (frontend) {Frontend\\Next.js};
    \node[comp, right=2.1cm of frontend] (backend) {Backend\\FastAPI};
    \node[data, below=1.3cm of backend] (postgres) {PostgreSQL};
    \node[infra, right=2.1cm of backend] (redis) {Redis\\Broker};
    \node[comp, right=2.1cm of redis] (worker) {Worker\\Celery};
    \node[data, below=1.3cm of worker] (qdrant) {Qdrant};
    \node[infra, above=1.2cm of worker] (ollama) {Ollama\\LLM + Embeddings};
    \node[obs, below=1.3cm of postgres] (prom) {Prometheus};
    \node[obs, right=2.1cm of prom] (grafana) {Grafana};

    \draw[line] (frontend) -- node[above]{HTTP REST} (backend);
    \draw[line] (backend) -- (postgres);
    \draw[line] (backend) -- (redis);
    \draw[line] (redis) -- (worker);
    \draw[line] (worker) -- (postgres);
    \draw[line] (worker) -- (qdrant);
    \draw[line] (backend) -- (qdrant);
    \draw[line] (backend) -- (ollama);
    \draw[line] (worker) -- (ollama);
    \draw[line] (prom) -- node[above]{scrape} (backend);
    \draw[line] (grafana) -- (prom);
  \end{tikzpicture}}
  \caption{Architecture logique globale}
\end{figure}

\section{Architecture de déploiement}
Le déploiement local et démonstratif repose sur Docker Compose. Chaque service est isolé dans un conteneur, avec des volumes dédiés pour PostgreSQL, Redis, Qdrant, Prometheus et Grafana. Le backend et le worker partagent un répertoire multiprocess Prometheus afin d'agréger les métriques issues des processus applicatifs.

\begin{table}[H]
  \centering
  \caption{Services de déploiement}
  \begin{tabularx}{\textwidth}{L{3cm}L{3cm}Y}
    \toprule
    \textbf{Service} & \textbf{Technologie} & \textbf{Responsabilité} \\
    \midrule
    frontend & Next.js & Interface utilisateur, authentification côté client, pages dashboard/documents/chat/settings. \\
    backend & FastAPI & API REST, sécurité, orchestration métier, métriques et documentation OpenAPI. \\
    celery\_worker & Celery & Traitement asynchrone des documents et indexation. \\
    postgres & PostgreSQL 16 & Données relationnelles, conversations, évaluations et métadonnées. \\
    redis & Redis 7 & Broker de tâches Celery. \\
    qdrant & Qdrant & Stockage et recherche vectorielle des chunks. \\
    prometheus & Prometheus & Collecte périodique des métriques. \\
    grafana & Grafana & Visualisation et tableaux de bord. \\
    \bottomrule
  \end{tabularx}
\end{table}

\section{Flux d'ingestion documentaire}
\begin{figure}[H]
  \centering
  \resizebox{\textwidth}{!}{
  \begin{tikzpicture}[
    step/.style={draw, rounded corners=5pt, fill=LightBlue, minimum width=2.8cm, minimum height=0.95cm, align=center},
    store/.style={draw, rounded corners=5pt, fill=LightTeal, minimum width=2.8cm, minimum height=0.95cm, align=center},
    line/.style={-Latex, thick, color=DevPilotGray}
  ]
    \node[step] (upload) {Upload\\PDF/TXT/MD};
    \node[step, right=1.2cm of upload] (validate) {Validation\\type + taille};
    \node[store, right=1.2cm of validate] (metadata) {Document\\queued};
    \node[step, right=1.2cm of metadata] (task) {Tâche\\Celery};
    \node[step, below=1.2cm of task] (extract) {Extraction\\texte};
    \node[step, left=1.2cm of extract] (clean) {Nettoyage\\chunking};
    \node[store, left=1.2cm of clean] (chunks) {Chunks\\PostgreSQL};
    \node[step, below=1.2cm of clean] (embed) {Embeddings\\Ollama};
    \node[store, right=1.2cm of embed] (qdrant) {Index\\Qdrant};
    \node[store, right=1.2cm of qdrant] (indexed) {Document\\indexed};

    \draw[line] (upload) -- (validate);
    \draw[line] (validate) -- (metadata);
    \draw[line] (metadata) -- (task);
    \draw[line] (task) -- (extract);
    \draw[line] (extract) -- (clean);
    \draw[line] (clean) -- (chunks);
    \draw[line] (clean) -- (embed);
    \draw[line] (embed) -- (qdrant);
    \draw[line] (qdrant) -- (indexed);
  \end{tikzpicture}}
  \caption{Pipeline d'ingestion et d'indexation}
\end{figure}

Le choix de Celery permet de dissocier le temps de réponse de l'API du temps de traitement documentaire. Le document est immédiatement persisté avec un statut \texttt{queued}, puis le worker extrait le texte, le nettoie, le découpe et indexe les vecteurs. Les erreurs sont capturées, les statuts sont mis à jour et les métriques de durée et d'échec sont enregistrées.

\section{Flux RAG agentique}
\begin{figure}[H]
  \centering
  \resizebox{\textwidth}{!}{
  \begin{tikzpicture}[
    agent/.style={draw, rounded corners=5pt, fill=LightOrange, minimum width=2.6cm, minimum height=0.9cm, align=center},
    data/.style={draw, rounded corners=5pt, fill=LightTeal, minimum width=2.6cm, minimum height=0.9cm, align=center},
    line/.style={-Latex, thick, color=DevPilotGray}
  ]
    \node[agent] (router) {Router};
    \node[agent, right=1cm of router] (rewrite) {Query\\Rewriter};
    \node[agent, right=1cm of rewrite] (retrieval) {Retrieval};
    \node[data, above=1.1cm of retrieval] (stores) {PostgreSQL\\+ Qdrant};
    \node[agent, right=1cm of retrieval] (rerank) {Reranker};
    \node[agent, right=1cm of rerank] (gen) {Generator};
    \node[agent, right=1cm of gen] (eval) {Evaluator};
    \node[agent, right=1cm of eval] (correct) {Corrector};
    \node[data, below=1.1cm of gen] (persist) {Messages\\Citations\\Traces};

    \draw[line] (router) -- (rewrite);
    \draw[line] (rewrite) -- (retrieval);
    \draw[line] (stores) -- (retrieval);
    \draw[line] (retrieval) -- (rerank);
    \draw[line] (rerank) -- (gen);
    \draw[line] (gen) -- (eval);
    \draw[line] (eval) -- (correct);
    \draw[line] (correct) |- (persist);
  \end{tikzpicture}}
  \caption{Workflow RAG agentique}
\end{figure}

Le workflow agentique améliore la maîtrise du comportement du RAG. Le routeur choisit une stratégie de recherche selon la question. Le réécrivain normalise la requête. Le retrieval interroge la recherche sémantique, lexicale ou hybride. Le reranker réordonne les chunks candidats. Le générateur produit une réponse avec le contexte. L'évaluateur attribue des scores de qualité. Le correcteur peut remplacer la réponse par une réponse d'insuffisance de contexte lorsque le risque d'hallucination est trop élevé.

\section{Choix de persistance}
PostgreSQL conserve les données transactionnelles : utilisateurs, workspaces, documents, chunks, conversations, messages, citations, évaluations, agent runs et usages LLM. Qdrant conserve les vecteurs et leurs payloads minimaux : workspace, document, chunk, index et nom du fichier. Ce découplage évite de forcer la base relationnelle à supporter la recherche vectorielle et permet d'utiliser chaque outil pour son domaine naturel.

\section{Scalabilité et résilience}
Le système est conçu pour évoluer par séparation des responsabilités :
\begin{itemize}
  \item le backend peut être répliqué horizontalement derrière un reverse proxy ;
  \item les workers Celery peuvent être multipliés selon le volume d'ingestion ;
  \item Redis absorbe le découplage entre upload et traitement ;
  \item Qdrant spécialise la recherche vectorielle ;
  \item Prometheus fournit les signaux nécessaires pour piloter la montée en charge.
\end{itemize}

\chapter{Réalisation technique}

\section{Backend FastAPI}
Le backend FastAPI expose des routes REST versionnées sous \texttt{/api/v1}. Les routes principales couvrent l'authentification, les workspaces, les documents, la recherche, le chat, la santé système et la configuration IA sûre. Les dépendances FastAPI assurent l'injection de session SQLAlchemy, l'authentification et la vérification d'accès aux workspaces.

La configuration est centralisée dans une classe Pydantic Settings. Les champs sensibles comme \texttt{SECRET\_KEY}, \texttt{DATABASE\_URL}, clés API ou tokens ne sont pas exposés par les endpoints publics de configuration. Les paramètres de RAG, embeddings, CORS, modèles et timeouts sont validés au démarrage.

\section{Frontend Next.js}
Le frontend est développé en Next.js avec TypeScript et Tailwind CSS. Il fournit des pages d'authentification, dashboard, documents, chat, workspaces et settings. Les composants réutilisables structurent l'état d'API, la navigation, les formulaires d'upload, la liste documentaire, les cartes de configuration et les retours d'erreur.

Le stockage du token JWT côté client permet l'accès aux routes protégées. Les pages fonctionnelles de documents et chat sont centrées sur le workspace, ce qui correspond au modèle de sécurité du backend.

\section{Providers IA}
Le système est prioritairement configuré pour Ollama en local. Le modèle de génération recommandé est \texttt{qwen3:8b}, tandis que les embeddings utilisent \texttt{nomic-embed-text}. Des points d'extension existent pour Gemini et OpenAI, principalement afin de démontrer une architecture provider-based.

\begin{table}[H]
  \centering
  \caption{Choix techniques majeurs}
  \begin{tabularx}{\textwidth}{L{3cm}L{3.2cm}Y}
    \toprule
    \textbf{Domaine} & \textbf{Choix} & \textbf{Justification} \\
    \midrule
    API & FastAPI & Performance, typage, OpenAPI natif, intégration Pydantic. \\
    Frontend & Next.js & Routage moderne, TypeScript, intégration React, productivité. \\
    ORM & SQLAlchemy async & Contrôle du modèle relationnel et compatibilité PostgreSQL. \\
    Migrations & Alembic & Versionnement du schéma et reproductibilité. \\
    Queue & Redis + Celery & Découplage des traitements longs et retries. \\
    Vector store & Qdrant & Recherche vectorielle dédiée, filtres par workspace, payloads. \\
    IA locale & Ollama & Développement privé, faible dépendance cloud, démonstration locale. \\
    Monitoring & Prometheus/Grafana & Standard industriel pour métriques et dashboards. \\
    Packaging & Docker Compose & Environnement reproductible pour démonstration et validation. \\
    \bottomrule
  \end{tabularx}
\end{table}

\section{Recherche hybride et reranking}
La recherche sémantique interroge Qdrant à partir de l'embedding de la requête. La recherche keyword repose sur PostgreSQL full-text search. La recherche hybride fusionne les résultats par Reciprocal Rank Fusion. Le reranking applique ensuite une heuristique de recouvrement lexical, de correspondances exactes, de score d'origine et de priorité de rang. Un reranker LLM optionnel peut être activé.

\section{Évaluation de la réponse}
L'évaluation automatique retourne quatre scores :
\begin{itemize}
  \item \textbf{faithfulness} : degré de support par le contexte ;
  \item \textbf{relevance} : adéquation à la question ;
  \item \textbf{context precision} : utilité des chunks récupérés ;
  \item \textbf{hallucination score} : risque de contenu non supporté.
\end{itemize}

Le mécanisme essaie d'abord une évaluation par LLM, puis bascule vers un fallback heuristique si le service IA n'est pas disponible ou si le JSON retourné est invalide. Cette décision est importante pour la résilience : l'indisponibilité partielle du modèle ne bloque pas toute la chaîne qualité.

\chapter{Sécurité}

\section{Authentification et sessions}
L'authentification repose sur JWT signé avec HS256 et une clé secrète configurée par variable d'environnement. Les mots de passe sont hachés avec bcrypt via Passlib. Le token contient le sujet utilisateur, l'email, le rôle et une expiration.

\section{Isolation multi-workspace}
L'isolation est une exigence centrale. Les opérations sur workspaces, documents, retrieval, chat et évaluation passent par des dépendances qui vérifient l'appartenance de l'utilisateur au workspace. L'utilisateur B ne peut pas accéder aux documents, recherches, chats ou évaluations de l'utilisateur A.

\begin{figure}[H]
  \centering
  \begin{tikzpicture}[
    zone/.style={draw, rounded corners=8pt, minimum width=5.4cm, minimum height=3.5cm, align=center},
    item/.style={draw, rounded corners=4pt, fill=LightGray, minimum width=3.8cm, minimum height=0.7cm, align=center},
    deny/.style={-Latex, thick, color=DevPilotRed},
    allow/.style={-Latex, thick, color=DevPilotGreen}
  ]
    \node[zone, fill=LightBlue] (a) {Workspace A};
    \node[item, below=0.25cm of a.north] (adoc) {Documents A};
    \node[item, below=0.2cm of adoc] (achat) {Chat A};

    \node[zone, fill=LightTeal, right=2.2cm of a] (b) {Workspace B};
    \node[item, below=0.25cm of b.north] (bdoc) {Documents B};
    \node[item, below=0.2cm of bdoc] (bchat) {Chat B};

    \node[above=0.8cm of a] (usera) {Utilisateur A};
    \node[above=0.8cm of b] (userb) {Utilisateur B};

    \draw[allow] (usera) -- (a);
    \draw[allow] (userb) -- (b);
    \draw[deny] (usera) to[bend left=15] node[above]{interdit} (b);
    \draw[deny] (userb) to[bend left=15] node[below]{interdit} (a);
  \end{tikzpicture}
  \caption{Principe d'isolation multi-workspace}
\end{figure}

\section{Validation des entrées}
Les schémas Pydantic imposent les contraintes de base : longueurs, UUID, types attendus et valeurs autorisées. L'upload documentaire applique une validation d'extension et une taille maximale. Le nom de fichier est nettoyé pour éviter les traversées de chemin. Le chemin de stockage est résolu et vérifié pour rester dans la racine de stockage autorisée.

\section{Gestion des secrets}
La configuration sépare les valeurs sensibles des informations exposables. Les endpoints de diagnostic renvoient uniquement des informations sûres : provider actif, modèles, paramètres de génération, activation du reranking, mais pas les clés API, mots de passe, JWT ou URL complètes sensibles.

\section{Analyse des risques}
\begin{table}[H]
  \centering
  \caption{Risques de sécurité et mesures}
  \begin{tabularx}{\textwidth}{L{3.2cm}Y Y}
    \toprule
    \textbf{Risque} & \textbf{Impact} & \textbf{Mesure} \\
    \midrule
    Vol ou fuite de token & Accès non autorisé & Expiration JWT, routes protégées, non-affichage des tokens dans l'UI. \\
    Accès cross-workspace & Fuite de documents privés & Vérification de membership sur chaque route métier. \\
    Upload malveillant & Écriture hors répertoire ou fichier non prévu & Nettoyage du nom, validation d'extension, contrôle de taille, résolution du chemin. \\
    Exposition de secrets & Compromission des providers et bases & Safe config, logs prudents, variables d'environnement. \\
    Hallucination IA & Mauvaise décision utilisateur & Citations, évaluation, correcteur, réponse d'insuffisance de contexte. \\
    Déni de service par gros fichier & Saturation CPU/mémoire & Limite de taille, chunking, traitement asynchrone. \\
    \bottomrule
  \end{tabularx}
\end{table}

\chapter{Tests et qualité logicielle}

\section{Stratégie de test}
La stratégie de test combine plusieurs niveaux :
\begin{itemize}
  \item tests unitaires backend pour les services, providers, modèles, sécurité, retrieval et agents ;
  \item tests d'intégration live pour valider la stack Docker, l'API, la base, l'authentification, les workspaces et les workflows ;
  \item tests d'ingestion avec Celery, Ollama et Qdrant ;
  \item tests RAG pour chat, citations, évaluation, agent runs et reranking ;
  \item smoke tests frontend pour vérifier le rendu des pages critiques.
\end{itemize}

\section{Résultats de validation}
Les résultats de validation documentés montrent une progression solide :
\begin{table}[H]
  \centering
  \caption{Synthèse des validations automatisées}
  \begin{tabularx}{\textwidth}{L{4cm}L{3cm}Y}
    \toprule
    \textbf{Suite} & \textbf{Résultat} & \textbf{Couverture} \\
    \midrule
    Intégration live sans ingestion/RAG & 19 passed, 10 deselected, 3 xfailed & Stack, santé, configuration, auth, workspaces, sécurité, routes frontend. \\
    Ingestion et retrieval live & 7 passed, 25 deselected & Upload, worker Celery, extraction, chunking, embeddings, Qdrant, retrieval. \\
    RAG, évaluation, agents, reranking & 3 passed, 29 deselected & Chat RAG, citations, évaluation persistée, agent runs, reranking. \\
    Compileall tests d'intégration & passed & Vérification syntaxique des tests. \\
    Tests unitaires backend host & bloqué environnement & Dépendances Python de développement non installées sur l'hôte. \\
    \bottomrule
  \end{tabularx}
\end{table}

\section{Critères d'acceptation}
\begin{table}[H]
  \centering
  \caption{État des critères d'acceptation principaux}
  \begin{tabularx}{\textwidth}{L{5cm}L{2cm}Y}
    \toprule
    \textbf{Critère} & \textbf{État} & \textbf{Observation} \\
    \midrule
    Services Docker démarrent correctement & \cmark & PostgreSQL, Redis, Qdrant, backend, worker et frontend validés. \\
    Authentification JWT fonctionnelle & \cmark & Register, login, me, erreurs négatives. \\
    Isolation workspace & \cmark & Accès croisés refusés dans les scénarios de sécurité. \\
    Upload et ingestion documentaire & \cmark & Statuts et indexation validés avec dépendances live. \\
    Retrieval sémantique, keyword et hybride & \cmark & Stratégies testées via API. \\
    Chat RAG avec citations & \cmark & Réponse, citations et persistance validées. \\
    Évaluation et correcteur & \cmark & Scores retournés et persistés. \\
    Observabilité Prometheus/Grafana & \cmark & Métriques et dashboard provisionnés. \\
    Route frontend globale documents & partiel & L'expérience fonctionnelle est workspace-aware. \\
    Route frontend globale chat & partiel & L'expérience fonctionnelle est workspace-aware. \\
    Sélection explicite de stratégie chat côté API & partiel & La stratégie est actuellement choisie par le routeur agentique. \\
    \bottomrule
  \end{tabularx}
\end{table}

\section{Qualité logicielle}
La qualité logicielle est renforcée par le typage, les schémas Pydantic, l'ORM, les migrations, les tests et la modularité. Le backend évite de mélanger contrôleurs HTTP et logique métier. Le code distingue erreurs métier, erreurs provider et erreurs d'infrastructure. Les métriques permettent de détecter les régressions de latence ou d'échec.

\chapter{Monitoring et observabilité}

\section{Objectifs}
L'observabilité vise à répondre rapidement à trois questions :
\begin{itemize}
  \item le système est-il disponible ;
  \item quelle partie du pipeline est lente ou en échec ;
  \item une réponse RAG peut-elle être reliée à sa requête, ses citations, ses agents et ses métriques.
\end{itemize}

\section{Logs structurés}
Chaque requête backend émet un événement de log \texttt{HTTP request completed}. Le log inclut un request id, la méthode, le chemin de route, le code HTTP, la latence, l'utilisateur si disponible et le workspace si applicable. Le backend renvoie aussi l'en-tête \texttt{X-Request-ID}, ce qui facilite la corrélation client/API/logs.

\section{Métriques Prometheus}
\begin{table}[H]
  \centering
  \caption{Métriques applicatives principales}
  \begin{tabularx}{\textwidth}{L{4.2cm}Y}
    \toprule
    \textbf{Métrique} & \textbf{Interprétation} \\
    \midrule
    \texttt{http\_requests\_total} & Volume de requêtes par méthode, route et statut. \\
    \texttt{http\_request\_duration\_seconds} & Latence HTTP par méthode, route et statut. \\
    \texttt{documents\_uploaded\_total} & Nombre de documents acceptés par type. \\
    \texttt{document\_processing\_duration\_seconds} & Durée des traitements Celery selon le statut final. \\
    \texttt{document\_processing\_failures\_total} & Échecs de traitement par type d'erreur. \\
    \texttt{rag\_queries\_total} & Nombre de requêtes RAG réussies ou échouées. \\
    \texttt{rag\_query\_duration\_seconds} & Durée de bout en bout du chat RAG. \\
    \texttt{retrieval\_latency\_seconds} & Latence du retrieval par stratégie et statut. \\
    \texttt{evaluation\_latency\_seconds} & Latence de l'évaluation LLM, heuristique ou erreur. \\
    \texttt{llm\_tokens\_total} & Tokens LLM par provider, modèle et type. \\
    \bottomrule
  \end{tabularx}
\end{table}

\section{Prometheus et Grafana}
Prometheus scrappe le backend sur \texttt{/metrics}. Grafana est provisionné avec une datasource Prometheus et un dashboard \og DevPilot AI Observability \fg{}. Cette organisation fournit une base professionnelle pour suivre le trafic, la latence, les erreurs, les traitements documentaires, la consommation LLM et la santé du retrieval.

\placeholder{Dashboard Grafana DevPilot AI Observability}{4cm}
\placeholder{Endpoint Prometheus avec métriques backend}{4cm}

\chapter{DevOps et exploitation}

\section{Environnement Docker Compose}
Docker Compose fournit un environnement reproductible. Les services disposent de healthchecks lorsqu'ils sont critiques, de volumes persistants et de ports configurables via variables d'environnement. L'application peut être lancée avec :

\begin{lstlisting}[language=bash]
docker compose up --build
\end{lstlisting}

Les ports peuvent être adaptés lorsque la machine hôte a déjà des services actifs :

\begin{lstlisting}[language=bash]
BACKEND_HOST_PORT=8001 FRONTEND_HOST_PORT=3001 \
FRONTEND_URL=http://localhost:3001 \
BACKEND_CORS_ORIGINS=http://localhost:3001 \
docker compose up --build
\end{lstlisting}

\section{Commandes d'ingénierie}
Le Makefile formalise les opérations fréquentes :
\begin{lstlisting}[language=bash]
make up
make down
make logs
make backend-install-dev
make test-backend
make test-integration
make test-ingestion
make test-rag
make test-cov
make migrate
make revision message="describe schema change"
\end{lstlisting}

\section{Migrations}
Alembic versionne le schéma PostgreSQL. Les migrations couvrent le schéma initial puis les évolutions liées à l'évaluation des réponses. Cette pratique est indispensable pour passer d'un prototype à un système maintenable, car elle rend l'état de la base reproductible.

\section{Configuration par environnement}
La configuration utilise des variables d'environnement et un fichier \texttt{.env} local ignoré par Git. Les valeurs exposées dans \texttt{.env.example} documentent les paramètres nécessaires sans publier de secrets réels. Les paramètres incluent les connexions PostgreSQL/Redis/Qdrant, les providers IA, les modèles, les limites RAG, CORS, uploads et métriques.

\section{Préparation production-ready}
Le projet présente plusieurs éléments compatibles avec une trajectoire production-ready :
\begin{itemize}
  \item conteneurisation des services ;
  \item configuration explicite et validée ;
  \item migrations de base de données ;
  \item séparation backend, worker, frontend et monitoring ;
  \item métriques et logs corrélables ;
  \item tests automatisés ;
  \item gestion d'erreurs claire ;
  \item isolation multi-workspace ;
  \item abstraction des providers IA.
\end{itemize}

Des étapes resteraient nécessaires pour une production stricte : TLS, reverse proxy, rotation des secrets, sauvegardes, politique de rétention, CI/CD complète, registry d'images, supervision d'alertes, durcissement réseau et gestion fine des rôles.

\chapter{Résultats}

\section{Résultat fonctionnel}
DevPilot AI permet de réaliser un scénario complet :
\begin{enumerate}
  \item démarrer la stack ;
  \item créer un utilisateur ;
  \item se connecter ;
  \item créer ou ouvrir un workspace ;
  \item importer un document ;
  \item attendre son indexation ;
  \item poser une question ;
  \item obtenir une réponse citée ;
  \item consulter l'évaluation ;
  \item suivre les traces et métriques.
\end{enumerate}

\section{Résultat architectural}
Le projet démontre une architecture cohérente de bout en bout. Le backend n'est pas un simple proxy vers un LLM : il orchestre des composants spécialisés, maintient un modèle de domaine, protège les accès, stocke les traces, expose des métriques et s'intègre avec un frontend.

\section{Résultat qualité}
Les validations montrent que les fonctionnalités critiques sont testables et majoritairement validées dans un environnement live. Les limites documentées sont explicitement identifiées, ce qui est un signe de maturité : un projet d'ingénierie crédible ne masque pas sa dette, il la qualifie.

\chapter{Limites et perspectives}

\section{Limites}
Les limites principales observées sont les suivantes :
\begin{itemize}
  \item l'expérience documents et chat est pleinement fonctionnelle dans le contexte workspace, tandis que certaines routes globales restent à aligner avec ce modèle ;
  \item la sélection explicite de stratégie de retrieval pour le chat n'est pas encore exposée dans le contrat d'API, car le routeur agentique choisit la stratégie ;
  \item l'environnement local dépend de la disponibilité d'Ollama et des modèles installés ;
  \item l'évaluation LLM peut basculer sur heuristique, ce qui améliore la disponibilité mais limite la finesse du jugement ;
  \item le projet ne comporte pas encore de CI/CD complète, d'alerting ou de durcissement cloud.
\end{itemize}

\section{Perspectives}
Les perspectives naturelles sont :
\begin{itemize}
  \item améliorer l'expérience workspace-first dans toutes les routes frontend ;
  \item exposer une sélection contrôlée des stratégies de retrieval ;
  \item renforcer le RAG par des stratégies d'évaluation plus avancées ;
  \item intégrer l'analyse automatique de code et les workflows DevOps ;
  \item ajouter un système d'alerting Prometheus ;
  \item mettre en place une CI/CD avec tests, linting, build d'images et déploiement ;
  \item enrichir les rôles et permissions au niveau workspace ;
  \item ajouter la gestion de collections par workspace ou par projet pour affiner l'isolation vectorielle.
\end{itemize}

\chapter{Conclusion}
DevPilot AI illustre la transformation d'une idée de chatbot documentaire en plateforme d'ingénierie complète. Le projet articule des choix modernes : FastAPI, Next.js, PostgreSQL, Redis, Celery, Qdrant, Ollama, Prometheus et Grafana. Il met surtout en évidence une pensée système : isoler les responsabilités, contrôler les données, tracer les réponses, évaluer la qualité, rendre le pipeline observable et préparer l'exploitation.

La valeur du projet réside dans sa cohérence globale. L'utilisateur ne reçoit pas simplement une réponse générée ; il obtient une réponse contextualisée, citée, évaluée et persistée. L'équipe technique dispose d'API testables, de métriques, de logs, de workers, de migrations et de commandes reproductibles. Cette combinaison donne à DevPilot AI une crédibilité supérieure à celle d'un prototype isolé et montre une démarche proche des attentes d'un environnement professionnel.

\appendix

\chapter{Annexe A -- Endpoints principaux}
\begin{longtable}{L{4.5cm}L{2.4cm}Y}
  \caption{API DevPilot AI}\\
  \toprule
  \textbf{Endpoint} & \textbf{Méthode} & \textbf{Rôle} \\
  \midrule
  \endfirsthead
  \toprule
  \textbf{Endpoint} & \textbf{Méthode} & \textbf{Rôle} \\
  \midrule
  \endhead
  \texttt{/health} & GET & Vérification simple de santé backend. \\
  \texttt{/metrics} & GET & Exposition des métriques Prometheus. \\
  \texttt{/api/v1/auth/register} & POST & Création d'un utilisateur. \\
  \texttt{/api/v1/auth/login} & POST & Authentification et émission de JWT. \\
  \texttt{/api/v1/auth/me} & GET & Profil utilisateur courant. \\
  \texttt{/api/v1/workspaces} & GET/POST & Liste et création de workspaces. \\
  \texttt{/api/v1/workspaces/\{id\}} & GET/PATCH & Consultation et modification d'un workspace. \\
  \texttt{/api/v1/workspaces/\{id\}/documents} & GET/POST & Liste et upload de documents dans un workspace. \\
  \texttt{/api/v1/documents/\{id\}/status} & GET & Statut de traitement d'un document. \\
  \texttt{/api/v1/workspaces/\{id\}/retrieval/search} & POST & Recherche sémantique, keyword ou hybride. \\
  \texttt{/api/v1/workspaces/\{id\}/chat/query} & POST & Question RAG dans un workspace. \\
  \texttt{/api/v1/messages/\{id\}/evaluation} & GET & Évaluation associée à un message. \\
  \texttt{/api/v1/system/ai-config} & GET & Configuration IA non sensible. \\
  \texttt{/api/v1/system/ai-health} & GET & Santé des providers IA. \\
  \bottomrule
\end{longtable}

\chapter{Annexe B -- Commandes de validation}
\begin{lstlisting}[language=bash]
docker compose config
docker compose up -d postgres redis qdrant backend celery_worker frontend
curl http://localhost:8000/health
curl http://localhost:8000/api/v1/system/ai-config

make backend-install-dev
make test-backend
make test-integration
DEVPILOT_RUN_INGESTION=1 make test-ingestion
DEVPILOT_RUN_INGESTION=1 DEVPILOT_RUN_RAG=1 make test-rag

docker compose up -d prometheus grafana
curl -sS http://localhost:8000/metrics
\end{lstlisting}

\chapter{Annexe C -- Configuration IA type}
\begin{lstlisting}
LLM_PROVIDER=ollama
EMBEDDING_PROVIDER=ollama
OLLAMA_BASE_URL=http://host.docker.internal:11434
OLLAMA_GENERATION_MODEL=qwen3:8b
OLLAMA_CHAT_THINK=false
OLLAMA_EMBEDDING_MODEL=nomic-embed-text
EMBEDDING_DIMENSION=768
GENERATION_TEMPERATURE=0.2
GENERATION_MAX_TOKENS=1000
RAG_TOP_K=5
RAG_MAX_CONTEXT_CHARS=6000
ENABLE_RERANKING=true
RETRIEVAL_CANDIDATES=15
RERANK_TOP_K=5
\end{lstlisting}

\chapter{Annexe D -- Placeholders de captures}
\placeholder{Page de connexion et inscription}{3.8cm}
\placeholder{Dashboard principal avec statut API}{3.8cm}
\placeholder{Liste de workspaces}{3.8cm}
\placeholder{Upload documentaire workspace-aware}{3.8cm}
\placeholder{Chat RAG avec citations et évaluation}{3.8cm}
\placeholder{Settings : configuration IA non sensible}{3.8cm}
\placeholder{Logs backend avec request id}{3.8cm}
\placeholder{Table PostgreSQL des agent runs}{3.8cm}

\begin{thebibliography}{99}

\bibitem{fielding2000}
R. T. Fielding, \textit{Architectural Styles and the Design of Network-based Software Architectures}, Doctoral dissertation, University of California, Irvine, 2000.

\bibitem{newman2015}
S. Newman, \textit{Building Microservices: Designing Fine-Grained Systems}, O'Reilly Media, 2015.

\bibitem{kleppmann2017}
M. Kleppmann, \textit{Designing Data-Intensive Applications}, O'Reilly Media, 2017.

\bibitem{owasp2021}
OWASP Foundation, \textit{OWASP Top 10: The Ten Most Critical Web Application Security Risks}, 2021.

\bibitem{nist80053}
National Institute of Standards and Technology, \textit{Security and Privacy Controls for Information Systems and Organizations}, SP 800-53 Rev. 5, 2020.

\bibitem{fastapi}
FastAPI Documentation, \textit{FastAPI framework, high performance, easy to learn, fast to code}, \url{https://fastapi.tiangolo.com/}.

\bibitem{sqlalchemy}
SQLAlchemy Documentation, \textit{SQLAlchemy Unified Tutorial and ORM Documentation}, \url{https://docs.sqlalchemy.org/}.

\bibitem{celery}
Celery Project, \textit{Distributed Task Queue}, \url{https://docs.celeryq.dev/}.

\bibitem{qdrant}
Qdrant Documentation, \textit{Vector Search Engine and Vector Database}, \url{https://qdrant.tech/documentation/}.

\bibitem{prometheus}
Prometheus Documentation, \textit{Monitoring system and time series database}, \url{https://prometheus.io/docs/}.

\bibitem{grafana}
Grafana Documentation, \textit{Grafana Observability Platform}, \url{https://grafana.com/docs/}.

\bibitem{docker}
Docker Documentation, \textit{Docker Compose Overview}, \url{https://docs.docker.com/compose/}.

\bibitem{rag}
P. Lewis et al., \textit{Retrieval-Augmented Generation for Knowledge-Intensive NLP Tasks}, Advances in Neural Information Processing Systems, 2020.

\end{thebibliography}

\end{document}
